\chapter{曲率}
在微分几何中，理解曲率的关键在于分清你是在观察\textbf{“空间中的线”}还是\textbf{“空间中的面”}。我们可以将这些曲率分为三个维度：曲线的几何量、曲面的分量、以及曲面的整体内蕴量。
\section{1. 曲线曲率 $\kappa$ 与曲率向量 $\vec{k}$}
这是最基础的量，描述的是一条特定的曲线在三维空间（或二维平面）中离开直线的程度。

曲线曲率 $\kappa$ 是一个一维对象（线）在$n$ 维空间中弯曲程度的通用度量。
平面的情况： 当曲线位于 $\mathbb{R}^2$ 平面时，它依然有曲率。此时，我们可以把平面看作是三维空间中 $z=0$ 的特殊子集。
其几何意义为： 在平面上，曲率 $\kappa$ 描述的是切线斜率（倾角 $\theta$）相对于弧长 $s$ 的变化率：
$$\kappa = \left| \frac{d\theta}{ds} \right|$$

在二维平面里，曲率可以有符号（左转为正，右转为负），因为法方向只有两个对称的选择。但在三维空间里，曲线可以向四面八方弯曲（还有“挠率”让它跳出平面），所以三维曲率 $\kappa$ 通常定义为非负的模长。

\begin{itemize}
\item \textbf{定义：} 设曲线按弧长参数 $s$ 参数化，其单位切向量为 $\vec{T}(s) = \frac{d\vec{r}}{ds}$。由于是弧长参数，速率恒定为 $1$，即内积满足 $\vec{T} \cdot \vec{T} = 1$--- 数学上的“恒定模长”约束。
\item \textbf{正交性推导：} 
根据向量内积的求导法则，对于随参数 $s$ 变化的向量 $\vec{u}(s)$ 和 $\vec{v}(s)$，有：$$ \frac{d}{ds}(\vec{u} \cdot \vec{v}) = \frac{d\vec{u}}{ds} \cdot \vec{v} + \vec{u} \cdot \frac{d\vec{v}}{ds} $$

于是，对等式 $\vec{T} \cdot \vec{T} = 1$ 两边关于 $s$ 求导，得到 $2 \vec{T} \cdot \frac{d\vec{T}}{ds} = 0$。这说明切向量的变化率 $\frac{d\vec{T}}{ds}$ 始终与切向量 $\vec{T}$ \textbf{垂直}（正交）。既然垂直，它就只能位于法方向 $\vec{n}$ 上。

\item \textbf{为何用法向刻画？} 离开直线的“快慢”之所以不用切向量而用法向刻画，是因为切向量的变化如果顺着切向，改变的是速度的大小（快慢）；只有变化量顺着法向，改变的才是轨迹的方向（形状）。因此，曲率向量定义为：
\[ \vec{k} = \frac{d\vec{T}}{ds} = \frac{d^2\vec{r}}{ds^2} = \kappa \vec{n} \]
其中 $\kappa$ 是标量曲率，$\vec{n}$ 是曲线的主法向量。该定义对二维平面曲线（$\kappa$ 可带符号）和三维空间曲线（$\kappa \geq 0$）同样适用。

\item \textbf{特点：} 它属于曲线自己，与这条线是否在某个曲面上无关。
\end{itemize}
\section{2. 法曲率 $\kappa_n$ (Normal Curvature)}
当你把一条曲线“贴”在曲面上时，曲线自身的弯曲 $\vec{k}$ 就被分解为曲面的特征分量。
\begin{itemize}
\item \textbf{定义：}法曲率定义为曲线曲率向量 $\vec{k}$ 在曲面单位法向量 $\vec{n}_{surf}$ 上的投影：

$$ \kappa_n = \vec{k} \cdot \vec{n}_{surf} = \frac{II}{I} $$

这里的“点乘”代表向量内积。
\item \textbf{几何直观：} 想象你沿着曲面某个方向切了一刀，切口（法截线）的弯曲程度就是法曲率。在物理直觉上，它代表了曲面在该方向上“上下起伏”的程度。

\item \textbf{联系：} 对于曲面上过同一点、同切方向的所有曲线，尽管它们的总曲率 $\kappa$ 可能不同，但它们在曲面法向上的投影（法曲率）都是一样的。
\end{itemize}
\section{3. 高斯曲率 $K$ (Gaussian Curvature)}
这是最深刻的量，它描述的是曲面本身的性质，不随观察角度或摆放方式改变。
\begin{itemize}
\item \textbf{定义：} 它是两个主曲率（法曲率在该点所有方向中的最大值 $\kappa_1$ 和最小值 $\kappa_2$）的乘积：

$$ K = \kappa\_1 \kappa\_2 $$
\item \textbf{绝妙之处：} 它是内蕴的。虽然定义里用了法曲率（外在投影量），但高斯证明了只需测量曲面上的距离（第一基本形式）就能算出来。
\end{itemize}
\section{核心区别与联系表}
\begin{table}[htbp]
\centering
\renewcommand{\arraystretch}{1.5}
\caption{曲线曲率、法曲率与高斯曲率对比}
\begin{tabular}{@{}llll@{}}
\toprule
\textbf{名称} & \textbf{描述对象} & \textbf{属性} & \textbf{直观含义} \\ \midrule
曲线曲率 $\kappa$ & 空间/平面线 & 外在（属于线） & 曲线偏离直线的程度（由 $\vec{T}$ 的转动率决定） \\
法曲率 $\kappa_n$ & 曲面上的方向 & 外在（属于嵌入） & 曲面在特定方向向“内/外”弯的程度 \\
高斯曲率 $K$ & 整个曲面局部 & \textbf{内蕴}（属于面本身） & 曲面局部像球($K>0$)、平地($K=0$)还是马鞍($K<0$) \\\bottomrule
\end{tabular}
\end{table}
\section{它们是如何串联起来的？}
为了理清逻辑，我们可以建立这样一个“投影链条”：
\begin{enumerate}
\item \textbf{第一步（运动产生）：} 在曲面上画一条线，由于切向量模长恒定，其导数产生垂直于切向的\textbf{总曲率向量 $\vec{k}$}。
\item \textbf{第二步（空间分解）：} 将 $\vec{k}$ 投影到曲面法向 $\vec{n}_{surf}$ 上，提取出属于地形特征的弯曲分量——\textbf{法曲率 $\kappa_n$}。
\item \textbf{第三步（极值筛选）：} 旋转切方向，寻找 $\kappa_n$ 的极大值和极小值，即\textbf{主曲率 $\kappa_1, \kappa_2$}。
\item \textbf{第四步（本质提取）：} 考察这两个极值的乘积，得到不随曲面在空间中如何弯折而改变的本质量——\textbf{高斯曲率 $K$}。
\end{enumerate}
\paragraph{总结：}
曲线曲率是“全貌”，法曲率是“侧影”，而高斯曲率是“本质”。
\paragraph{核心理解补充：}
\begin{itemize}
\item \textbf{曲线曲率 $\vec{k}$：} 就像是汽车行驶轨迹的总偏转，无论地面平坦还是坎坷，它只记录轨迹方向的变化。
\item \textbf{法曲率 $\kappa_n$：} 则是将这种偏转投影到地面的垂直方向上。这好比在山路上开车，法曲率告诉你的是这段路“由于地势原因”带来的上下起伏感。
\item \textbf{高斯曲率 $K$：} 最终定义了地面的形状。如果两个主曲率同号（如球体），意味着地面向同一个方向弯曲（同为凸或同为凹）；如果异号（如马鞍面），则意味着一个方向向上凹，另一个方向向下凸。
\end{itemize}
